
\section{\trGenus{} $0$ 的\trCurve}

现在，我们已经准备好（以某种形式）回答这样一个问题了：\trGenus{} $0$ 的\trCurve{}都是些什么？

在 \S7.5.10 中，我们见过一条（域 $k$ 上的）\trGenus{} $0$ \trCurve，它并不同构于 $\mathbb{P}^1$：即 $\mathbb{P}^2_\mathbb{R}$ 中的 $x^2+y^2+z^2=0$。（由 \eqref{eq:genus-plane-curve}，它的\trGenus{}为 $0$。）我们已经观察到，这条曲线并不同构于 $\mathbb{P}^1_\mathbb{R}$，因为它没有 $\mathbb{R}$-\trValuedPoint。另一方面，我们还不曾在\trAlgClosedField{}上见过具有这种性质的\trGenus{} $0$ \trCurve。这绝非巧合：$k$-\trValuedPoint{}的缺失，正是一条\trGenus{} $0$ \trCurve{}成为 $\mathbb{P}^1$ 的唯一障碍。

\textbf{19.3.1.\ \trhProposition。}\ ——\ \emph{设 $C$ 的\trGenus{}为 $0$，且 $C$ 有一个 $k$-\trValuedPoint（\trDegree{} $1$）。那么 $C\cong\mathbb{P}^1_k$。}

于是我们看到：\trAlgClosedField{}上的一切\trGenus{} $0$（\trIntegral、\trRegular、\trProjective）\trCurve{}都同构于 $\mathbb{P}^1$。

\emph{证明。}\quad 设 $p$ 为那个点，并考虑 $\mathscr{L}=\OO_C(p)$。于是 $\deg\mathscr{L}=1$，故我们可以应用上文已知的结果：其一，$h^0(C,\mathscr{L})=2$（\mbox{注记~19.2.5}）；其二，这两个\trSection{}给出一个到 $\mathbb{P}^1_k$ 的\trClosedEmbedding（\mbox{注记~19.2.11}，经\mbox{练习~19.2.E} 扩充）。但是，一条\trCurve{}到\trIntegral\trCurve{} $\mathbb{P}^1_k$ 之内的\trClosedEmbedding{}只能是同构！\hfill$\square$

附带一提，\mbox{命题~19.3.1} 蕴涵：$\mathbb{P}^2_\mathbb{R}$ 中的 $x^2+y^2+z^2=0$ 在 $\mathbb{R}$ 上没有\emph{次数 $1$ 的\trLineBundle}——否则，我们只须把上面的论证施加到相应的\trLineBundle{}上就行了。这个例子向我们表明：在一个非\trAlgClosed{}的域上，可以存在不同构于 $\mathbb{P}^1_k$ 的\trGenus{} $0$ \trCurve。下面这个结果让我们也能把它们亲手抓到。

\textbf{19.3.2.\ \trhClaim。}\ ——\ \emph{一切\trGenus{} $0$ \trCurve{}都可以描述为 $\mathbb{P}^2_k$ 中的\trConic。}

\emph{证明。}\quad 任何\trGenus{} $0$ \trCurve{}都有一个\trDegree{}为 $-2$ 的\trLineBundle——即\trCanonical{}层 $\omega_C$。于是任何\trGenus{} $0$ \trCurve{}都有一个\trDegree{}为 $2$ 的\trLineBundle：$\mathscr{L}=\omega_C^{\vee}$。我们应用\mbox{练习~19.2.E}：$\deg\mathscr{L}=2\ge 2g+1$，故这个\trLineBundle{}给出一个\trClosedEmbedding{} $\iota\colon C\hookrightarrow\mathbb{P}^2$。\mbox{练习~18.6.H} 保证该像作为\trClosedSubscheme{}的\trDegree{}为 $2$；但我们还没有完全证明，这条\trCurve{}可以由单个\trDegree{} $2$ 的齐次多项式割出。（这件事比你想的更微妙——在读下去之前，你不妨先试着自己论证一番。）我们有同构 $\Sym^2\mathrm{H}^0(\mathbb{P}^2,\OO(1))\overset{\sim}{\longrightarrow}\mathrm{H}^0(\mathbb{P}^2,\OO(2))$（两边都是三个变量的齐次二次多项式）。考察映射 $\Sym^2\mathrm{H}^0(C,\mathscr{L})\to\mathrm{H}^0(C,\mathscr{L}^{\otimes 2})$，由 \eqref{eq:h0-high-degree}，它是从一个维数为 $6$ 的\trVectorSpace{}到一个维数为 $5$ 的\trVectorSpace{}的映射。这个映射有非平凡的\trKernel，故 $\iota(C)$ \trSchemeTheoretic{}地含于 $V(f)$ 之中，其中 $f$ 为某个齐次\trConic。

\textbf{19.3.A.（练习）}\quad 用\mbox{练习~18.6.E} 证明：$\iota(C)$ \trSchemeTheoretic{}地等于 $V(f)$，从而完成上述证明。

\begin{flushright}
$\square$
\end{flushright}

\textbf{19.3.B.（练习）}\quad 设 $C$ 是一条\trGenus{} $0$ \trCurve（\trProjective、\trGeometricallyIntegral，且\trRegular）。证明：$C$ 有一个\trDegree{}至多为 $2$ 的点。（点的\trDegree{}曾定义于 \S5.3.8。）

\S7.5.9 中所呈现的、用几何手段寻找勾股三元组的方法，看上去颇为不同，但其实是一回事。当时有一条\trGenus{} $0$ \trCurve{} $C$（一条平面\trConic），带有一个 $k$-\trValuedPoint{} $p$，而我们证明了它同构于 $\mathbb{P}^1_k$。用以证明该同构的\trLineBundle{}并不是\trDegree{} $1$ 的\trLineBundle{} $\OO_C(p)$；而是\trDegree{} $1$ 的\trLineBundle{} $\OO_{\mathbb{P}^2}(1)|_C\otimes\OO_C(-p)$。
